\scanpage{8}
\mspage{7}
\mspage{-0.14-}
\uline{Proof}: Note first of all that $\phi$ is dominating by (ii); hence by (iii), $\phi$ is actually \uline{surjective}. Now suppose $\deleted{\phi}\inserted{\psi}:X\longrightarrow Z$ is any $S$-morphism such that $\psi\circ\sigma=\psi\circ p_2$. We proceed exactly as in the proof of Proposition 0.1: choose an affine open covering $\{V_i\}$ of $Z$. Once we show that there is an open covering $\{U_i\}$ of $Y$ such that
\figure{\begin{tikzpicture}\node[draw,ellipse,inner sep=3pt] {$\psi^{-1}(V_i)\supset\phi^{-1}(U_i)$};\end{tikzpicture}},
we can conclude the proof as in Proposition 0.1. To show this, let $W_i=X-\psi^{-1}(V_i)$. These are closed invariants subsets of $X$. Therefore, by (iii), $U_i=Y-\phi W_i$ are open; and $\phi^{-1}(U_i)\subset\psi^{-1}(V_i)$ for all $i$. Finally, $\{\psi^{-1}(V_i)\}$ cover $X$, hence $\bigcap_iW_i=\phi$, hence by (iii) $\bigcap_i\phi W_i=\phi$, hence $\{U_i\}$ cover $Y$. \uline{QED}

(7). Suppose $(Y,\phi)$ is a geometric quotient of $X$ by $G$. Notice that conditions (i), (ii), and (iii) of Definition 0.6 are preserved by any base extension $Y'\to Y$. Only the more delicate condition (iv) might (and indeed sometimes does) break down. However, if $Y'\longrightarrow Y$ is \uline{flat}, then condition (iv) is preserved. To see this, let $\psi=\phi\circ\sigma=\phi\circ p_2$. Then (iv) asserts the exactness of:
\margin{$\varphi$ quasi-compact quasi-séparé}
\diagram{\begin{tikzcd}[ampersand replacement=\&,column sep=large]
(*)\quad 0\arrow[r] \& \underline o_Y\arrow[r] \& \phi_*(\underline o_X)\arrow[r,"\sigma^*-p_2^*"] \& \psi_*(\underline o_{G\times_S X})
\end{tikzcd}}
But taking direct image sheaves commutes with \uline{flat} base extension, i.e.
\[
\begin{aligned}
\phi_*(\underline o_X)\otimes\underline o_Y&\cong\phi'_*(\underline o_X)\\
\psi_*(\underline o_{G\times_S X})\otimes\underline o_Y&=\psi'_*(\underline o_{G\times_S X}).
\end{aligned}
\]
(where $X'=X\times_Y Y'$, etc.), and the exactness of a sequence is pre-
